% ---
% titre: Recette de crêpes (tableau de proportionnalité, quatrième proportionnelle)
% theme: FA
% chapitre: Proportionnalité
% sous_chapitre: Résoudre des problèmes de proportionnalité
% niveau: [10e]
% type: EVAL
% tags: [c-proportionnalite,p-question-TS]
% difficulte: 1
% corrige: true
% ---

\question
Une voiture consomme en moyenne 6,4 litres d'essence pour 100 kilomètres.
 
\begin{parts}
\vspace{20pt}\part[2] Combien de litres d'essence consomme-t-elle pour un trajet de 275\,km ?
 
\vspace{10pt}{\footnotesize\raggedleft Espace pour ta démarche et tes calculs\par}
\begin{solutionorgrid}[120pt]
 
\vspace{5pt}
\def\arraystretch{1.5}
\begin{tabularx}{0.5\textwidth}{l|C|C|}
Distance [km]&100&275\\
\hline
Essence [l]&6,4&\\
\end{tabularx}
\def\arraystretch{1}
 
\hfill\textbf{(0,5\,pt)} pour la construction du tableau
 
\vspace{15pt}$275 \cdot 6,4 : 100 = 17,6$\,l  \hfill calcul \textbf{(0,5\,pt)} et réponse \textbf{(0,5\,pt)}
\end{solutionorgrid}
 
\begin{tcolorbox}[enhanced jigsaw, interior hidden, colframe=box, parbox=false]%
Réponse : 
\sol{\vspace{20pt}}{Pour un trajet de 275\,km, elle consomme 17,6 litres d'essence. \textbf{(0,5\,pt)}}
\end{tcolorbox}
 
\vspace{20pt}\part[2] Le réservoir de la voiture contient 45 litres d'essence. Quelle distance peut-elle parcourir avec un plein complet?
 
\vspace{10pt}{\footnotesize\raggedleft Espace pour ta démarche et tes calculs\par}
\begin{solutionorgrid}[120pt]
 
\vspace{5pt}
\def\arraystretch{1.5}
\begin{tabularx}{0.5\textwidth}{l|C|C|}
Distance [km]&100&\\
\hline
Essence [l]&6,4&45\\
\end{tabularx}
\def\arraystretch{1}
 
\hfill\textbf{(0,5\,pt)} pour la construction du tableau
 
\vspace{15pt}$45 \cdot 100 : 6,4 = 703,125$\,km  \hfill calcul \textbf{(0,5\,pt)} et réponse \textbf{(0,5\,pt)}
\end{solutionorgrid}
 
\begin{tcolorbox}[enhanced jigsaw, interior hidden, colframe=box, parbox=false]%
Réponse : 
\sol{\vspace{20pt}}{Avec un plein complet, elle peut parcourir environ 703\,km. \textbf{(0,5\,pt)}}
\end{tcolorbox}
\end{parts}